% Spanish for Beginners — Week 3 Study Guide (Unit 2)
% Compile: pdflatex unit-2-por-para.tex

\documentclass[10pt,a4paper]{article}
\usepackage[utf8]{inputenc}
\usepackage[T1]{fontenc}
\usepackage[margin=1.2cm]{geometry}
\usepackage{enumitem}
\setlist{nosep,leftmargin=*}
\usepackage{booktabs}
\usepackage{xcolor}
\usepackage{fancyhdr}
\usepackage{microtype}
\definecolor{w3purple}{RGB}{80,60,120}
\pagestyle{fancy}
\fancyhf{}
\fancyfoot[L]{\small Spanish for Beginners | Week 3}
\fancyfoot[R]{\small\thepage}
\raggedright

\begin{document}
\begin{center}
{\large\bfseries\color{w3purple} Week 3 -- Unit 2}\\[0.2em]
\footnotesize Por/para/porque, actividades para aprender.
\end{center}\vspace{0.5em}

\section*{\color{w3purple}Por vs.\ Para vs.\ Porque}
\textbf{Para} + infinitivo = purpose (in order to). \textit{Quiero aprender espa\~nol para hablar con mis amigos.} \\
\textbf{Porque} + clause = reason (because). \textit{Quiero visitar el Prado porque quiero ver los cuadros.} \\
\textbf{Por} + sustantivo = reason (for/because of). \textit{Quiero visitar Roma por la comida y los monumentos.} \\
\textbf{Trick:} Para + infinitive. Porque + conjugated verb. Por + noun.

\section*{\color{w3purple}Actividades para aprender}
\textbf{ver:} pel\'iculas, series, la televisi\'on, v\'ideos. \textbf{escuchar:} m\'usica, la radio, podcasts. \\
\textbf{leer:} libros, novelas, art\'iculos. \textbf{escribir:} un diario, mensajes, emails. \\
\textbf{practicar:} con nativos, la pronunciaci\'on, la gram\'atica. \textbf{hacer:} ejercicios, deberes, un intercambio.

\section*{\color{w3purple}Lenguas por pa\'is}
Rusia $\to$ ruso. Italia $\to$ italiano. Francia $\to$ franc\'es. Alemania $\to$ alem\'an. \\
Brasil $\to$ portugu\'es. China $\to$ chino. Jap\'on $\to$ japon\'es. Corea $\to$ coreano.
\end{document}
